% ============================================================
% section5_knowledge.tex — Section V: Knowledge & Continuity
% ============================================================

\chapter{Knowledge \& Continuity}

\sectionintro{``The knowledge a frontier team generates is as valuable as
the artifacts it produces. Leaders who treat it as such build organizations
that learn. Leaders who don't build organizations that repeat themselves.''}

Frontier work generates hard-won knowledge at an extraordinary rate.
Most organizations waste most of it. These principles address the discipline
of capturing, transferring, and building on institutional knowledge---so that
each generation of the team starts from a higher floor than the last.

% ----------------------------------------------------------
\section{Apply Academic Research Rigor}

\principle{Research-grade documentation discipline is a competitive advantage in business R\&D.}

Academic research has solved the knowledge continuity problem in a specific way:
rigorous documentation of methods, results, and reasoning, structured so that
others can build on the work without having to reconstruct it from scratch.
Business R\&D rarely applies this discipline. It should.

Applying research rigor to R\&D documentation does not mean bureaucracy.
It means being deliberate about what is captured---not just what was built,
but why, what was tried and failed, what assumptions were made, and what
the open questions are. This is the difference between a project archive
and a knowledge asset.

\subsection{Story}
\tbd[Story to be attached. Note: the Red Program Tracing story is confirmed
at V.2 (Stand on Each Other's Shoulders) and is no longer a candidate here.
This principle needs its own story from Eric.]

% ----------------------------------------------------------
\section{Stand on Each Other's Shoulders}

\principle{Each generation of the team should start from a higher floor than the last.}

The goal of knowledge continuity is not preservation---it is elevation.
When knowledge is captured and transmitted effectively, each new team member,
each new program phase, each new generation of the work begins from a higher
starting point. The team doesn't just avoid repeating mistakes; it compounds
its understanding over time.

This is the R\&D equivalent of the scientific tradition of standing on the
shoulders of giants. Applied within an organization, it means that the
institutional knowledge accumulated by one cohort becomes the foundation
for the next.

\begin{storybox}
\textbf{The Red Program Tracing Story} \textit{(placement confirmed)}

The team traced every significant decision back through the program's living
documentation. New members didn't just learn what had been decided---they
could reconstruct the full reasoning behind each decision without asking
anyone. The team hadn't just done good record-keeping. They had built a
living knowledge system: a knowledge-transfer mechanism that worked without
its authors being in the room.

\tbd[Skeleton from the raw capture. Full narrative---program context,
specifics, outcome---to be developed from Eric's direct narration.]
\end{storybox}

% ----------------------------------------------------------
\section{Principle 3}

\principle{\textit{\color{midgray}[TBD: Principle title and thesis to be confirmed with Eric.]}}

\tbd[Placeholder for Section V Principle 3.]
